\section{Discussion and Limitations}
\label{sec:discussion}
The study locates two distinct obstacles. Early local corrections often
leave the policy facing an unresolved prerequisite or phase transition.
Making a larger part of the strategy executable can complete a task, but
that success may belong entirely to the intervention program. The later
dual-view update does not transfer this observed utility into higher
scaffold-free selection success.

The learner receives initial-action targets from just four source states,
whereas evaluation requires it to reach and complete these decisions from
task reset. This mismatch in both coverage and starting conditions offers
a plausible explanation for weak transfer. A controlled training comparison
is needed to establish its causal role. The trajectories
also reveal persistent menu and prerequisite errors. A useful next test
would match source states and training budget while varying coverage of
the verified continuation, before scaling the number of rounds. In particular, the current result
suggests measuring what the learner does after each verified intermediate
state, as well as from task reset. This would distinguish learning the
successful continuation from learning to reach its starting point.
Such a comparison follows directly from the observed gap; its outcome
remains an empirical question.

The revisions and portfolios were chosen sequentially, so their outcome
counts are descriptive. The 355-task comparison uses one seed and one
candidate; it measures neither training variability nor final benchmark
generalization. Matched action-only, filtering, equal-compute, and research-memory
comparisons, together with a final RePair benchmark evaluation, remain
open tests. The heterogeneous released baselines provide task context.
Our present evidence concerns the local-to-policy transfer gap in this
development study.

Five paired repetitions establish the registered stability criterion,
not statistical certainty, and failure-selected states provide limited
opportunities to measure harm. Human engineering recovery and a disclosed
promotion exception prevent interpreting the development campaign as
fully unattended. RePair's contribution here is the explicit separation
and measurement of these handoffs, including the observed absence of a retained policy gain.
